\appendix
% Give appendices letter-based hyperref destinations (chapter.A/B/C, section.A.1 ...)
% so bookmarks and \ref links land in the appendix instead of colliding with the
% \theHchapter of main Chapters 1--3. \theHsection inherits this automatically.
\renewcommand{\theHchapter}{\Alph{chapter}}
\chapter{Appendix A: Proofs and Definitions}

\section{The Wasserstein-1 distance and its dual}\label{app:wasserstein}

The Wasserstein-1 distance measures how far apart two distributions $P_r$ and $P_g$ are by the smallest average distance the mass of one must be moved to match the other. Writing $\Pi(P_r,P_g)$ for the couplings, the joint distributions of $(X,X')$ with marginals $P_r$ and $P_g$,
\begin{equation}
W_1(P_r,P_g)=\inf_{\gamma\in\Pi(P_r,P_g)}\mathbb{E}_{(X,X')\sim\gamma}\bigl[\,\|X-X'\|\,\bigr].
\label{eq:wasserstein-primal}
\end{equation}
Each coupling $\gamma$ is a transport plan and the expectation is its average moving cost, so the infimum is the cheapest plan. This form cannot be optimised during training because it ranges over all couplings. The Kantorovich--Rubinstein duality rewrites the same distance as a maximisation over functions,
\begin{equation}
W_1(P_r,P_g)=\sup_{\|f\|_L\le 1}\bigl\{\mathbb{E}_{X\sim P_r}[f(X)]-\mathbb{E}_{X'\sim P_g}[f(X')]\bigr\}.
\label{eq:wasserstein-dual}
\end{equation}
The function $f$ is required to be $1$-Lipschitz, $|f(x)-f(x')|\le\|x-x'\|$. That constraint is what makes the maximum equal the transport cost: because $f$ can differ between two points by at most their distance, the gap between its average values under the two distributions cannot exceed the cost of moving the mass, and the best $f$ attains that cost. The dual is the form a Wasserstein generative adversarial network uses, since $f$ can be approximated by a critic network whose output is a real-valued score. Without the Lipschitz constraint the supremum is unbounded whenever $P_r$ and $P_g$ differ, and the critic would then report an arbitrary score in place of $W_1(P_r,P_g)$.

\section{Raking and AIPW under quadratic calibration}\label{app:calib-aipw}

Let $\pi_i$ be phase-two inclusion probabilities, $w_i:=\pi_i^{-1}$ the base weights, $H_i$ fully observed auxiliaries (phase-one), and $U_i(\beta)$ the estimating function for the parameter of interest $\beta$. Denote the full-cohort auxiliary total by $T_H=\sum_{i=1}^N H_i$.

\begin{proposition}[Raking $=$ AIPW under quadratic calibration]\label{prop:raking-aipw}
Consider \emph{generalised raking} with the quadratic (GREG) distance $d_1(a,b)=(a-b)^2/(2b)$:
\[
\min_{\{w_i^{\rm cal}\}}\ \sum_{i=1}^N R_i\,d_1\!\big(w_i^{\rm cal},\,w_i\big) \quad\text{s.t.}\quad \sum_{i=1}^N R_i\,w_i^{\rm cal} H_i \;=\; T_H.
\]
Its solution satisfies $w_i^{\rm cal} = w_i\bigl(1+\lambda^\top H_i\bigr)$ for some Lagrange multiplier $\lambda$ (unique when $\sum R_i w_i H_iH_i^\top$ is nonsingular). Let $\widehat\Gamma(\beta)$ be the base-weighted least-squares projection of $U_i(\beta)$ onto $H_i$, defined by
\[
\sum_{i=1}^N R_i w_i\,[U_i(\beta)-\widehat\Gamma(\beta)H_i]H_i^\top=0 .
\]
Then the $\beta$ that solves the raked score $\sum_i R_i w_i^{\rm cal} U_i(\beta)=0$ also solves the AIPW estimating equation
\[
\sum_{i=1}^N\Big\{\, R_i w_i\big[U_i(\beta)-\widehat\Gamma(\beta)H_i\big] \;+\; \widehat\Gamma(\beta)\,H_i \Big\}\;=\;0.
\]
Thus, under quadratic calibration, the raked score equation and the AIPW estimating equation are algebraically identical as functions of $\beta$ on the same realised data.
\end{proposition}

\begin{proof}
By the calibration constraint, $\sum_i H_i=\sum_i R_i w_i^{\rm cal}H_i$. Write the AIPW left-hand side at an arbitrary $\beta$:
\[
\sum_i R_i w_i\big[U_i(\beta)-\widehat\Gamma H_i\big] + \widehat\Gamma\sum_i H_i \;=\; \sum_i R_i\Big\{ w_i\big[U_i(\beta)-\widehat\Gamma H_i\big] + w_i^{\rm cal}\,\widehat\Gamma H_i \Big\}.
\]
Let $A_{\mathrm{AIPW}}(\beta)$ and $A_{\mathrm{rak}}(\beta)$ denote, respectively, the AIPW left-hand side and the raked-score left-hand side. Insert and subtract $w_i^{\rm cal}U_i(\beta)$, then regroup:
\[
A_{\mathrm{AIPW}}(\beta)-A_{\mathrm{rak}}(\beta)
=-\sum_i R_i\big(w_i^{\rm cal}-w_i\big)\big[U_i(\beta)-\widehat\Gamma H_i\big].
\]
Under quadratic calibration $w_i^{\rm cal}-w_i = w_i\,\lambda^\top H_i$, so
\[
A_{\mathrm{AIPW}}(\beta)-A_{\mathrm{rak}}(\beta)
=-\Bigg\{\sum_i R_i w_i\,[U_i(\beta)-\widehat\Gamma H_i]H_i^\top\Bigg\}\lambda .
\]
The defining normal equation for the base-weighted projection sets the term in braces to zero, so $A_{\mathrm{AIPW}}(\beta)=A_{\mathrm{rak}}(\beta)$ for every $\beta$. The two estimating functions have the same finite-sample root and the same asymptotic distribution. This cancellation requires $\widehat\Gamma$ to be the base-weighted projection; the calibrated-weight projection would not set the term in braces to zero.
\end{proof}

\begin{remark}\label{rem:raking-distance}
The exact multiplicative form $w_i^{\rm cal}=w_i\bigl(1+\lambda^\top H_i\bigr)$ is specific to the quadratic (chi-square) distance $d_1$, whose stationarity condition is linear in $w_i^{\rm cal}$. The entropy and logit distances more commonly intended by the term \emph{generalised raking} reproduce this form only to first order in $\lambda$, and, unlike $d_1$, they constrain the calibrated weights to be positive. The equivalence with the augmented estimating equation is therefore exact for $d_1$ and only first-order for the other distances.
\end{remark}

\section{Fourth-root plug-in argument}\label{app:fourth-root}

\begin{proposition}[Fourth-root plug-in]\label{prop:fourth-root}
Let $U(\beta,\alpha)=\frac{1}{n}\sum_{i=1}^n U_i(\beta,\alpha)$ be an estimating equation for the parameter of interest $\beta$ with nuisance parameter $\alpha$. Suppose:
\begin{enumerate}\setlength\itemsep{2pt}
\item[(A1)] \emph{Identification, smoothness and empirical limits:} $\mathbb{E}\{U_i(\tilde\beta,\alpha^*)\}=0$, with $\tilde\beta$ the unique root of $\beta\mapsto\mathbb{E}\{U_i(\beta,\alpha^*)\}$ in a neighbourhood of $\tilde\beta$, and $A_\beta:=\mathbb{E}\{-\partial U_i(\tilde\beta,\alpha^*)/\partial\beta\}$ is nonsingular; $U_i(\beta,\alpha)$ is twice continuously differentiable on a neighbourhood $\mathcal{N}$ of $(\tilde\beta,\alpha^*)$ with an integrable envelope $\sup_{(\beta,\alpha)\in\mathcal{N}}\lVert\partial^2 U_i(\beta,\alpha)\rVert\le M_i$, $\mathbb{E}M_i<\infty$; differentiation may be interchanged with expectation in $\mathcal{N}$; the sample average $n^{-1}\sum_i U_i(\beta,\alpha)$ obeys a uniform law of large numbers over $\mathcal{N}$, jointly in $(\beta,\alpha)$; and the centred empirical averages of $U_i(\tilde\beta,\alpha^*)$ and $\partial_\alpha U_i(\tilde\beta,\alpha^*)$ are $O_p(n^{-1/2})$. These empirical conditions may follow from an independent and identically distributed central limit theorem or from the corresponding central limit theorem for the sampling design.
\item[(A2)] \emph{Orthogonality (Neyman):} $\partial_\alpha \mathbb{E}\{U_i(\tilde\beta,\alpha)\}\big|_{\alpha=\alpha^*}=0$.
\item[(A3)] \emph{Rate:} $\|\hat\alpha - \alpha^*\|_2 = o_p(n^{-1/4})$, where $\alpha^* = \mathrm{plim}\,\hat\alpha$.
\end{enumerate}
Then
\[
\sqrt{n}\,(\hat\beta(\hat\alpha) - \tilde\beta) \;=\; A_\beta^{-1}\Big(\frac{1}{\sqrt{n}}\sum_{i=1}^n U_i(\tilde\beta, \alpha^*)\Big) \;+\; o_p(1).
\]
In particular, the first-order asymptotic distribution of $\hat\beta(\hat\alpha)$ depends on $\hat\alpha$ only through its probability limit $\alpha^*$, so that the estimation noise in $\hat\alpha$ does not affect the limiting variance of $\hat\beta$.
\end{proposition}

\begin{proof}
Taylor-expand $U(\beta,\alpha)$ at $(\tilde\beta,\alpha^*)$:
\begin{align*}
0 = U(\hat\beta(\hat\alpha),\hat\alpha)
&= U(\tilde\beta, \alpha^*)
+ \partial_\beta U(\tilde\beta, \alpha^*)\,(\hat\beta(\hat\alpha) - \tilde\beta) \\
&\quad + \partial_\alpha U(\tilde\beta, \alpha^*)\,(\hat\alpha - \alpha^*) \\
&\quad + O_p\!\big(\|\hat\beta(\hat\alpha) - \tilde\beta\|^2 + \|\hat\alpha - \alpha^*\|^2\big).
\end{align*}
By (A1), $U(\tilde\beta, \alpha^*)=O_p(n^{-1/2})$. The uniform law of large numbers in (A1) gives $\partial_\beta U(\tilde\beta, \alpha^*) \to -A_\beta$ in probability. By (A2) and the interchange condition in (A1), $\mathbb{E}\{\partial_\alpha U(\tilde\beta,\alpha^*)\}=0$, and the empirical-limit condition in (A1) gives $\partial_\alpha U(\tilde\beta,\alpha^*)=O_p(n^{-1/2})$. Combined with (A3), the cross-term satisfies $\partial_\alpha U(\tilde\beta,\alpha^*)\,(\hat\alpha-\alpha^*)=o_p(n^{-1/2})$. Preliminary consistency, $\hat\beta(\hat\alpha)\xrightarrow{p}\tilde\beta$, follows from the local uniqueness of the root and the uniform law of large numbers assumed in (A1), together with (A3), by the standard argument for $M$-estimators with an estimated nuisance \citep{newey1994large}. Given consistency, the linearisation is self-consistent: the leading terms are $O_p(n^{-1/2})$ and the quadratic remainder $O_p(\lVert\hat\beta(\hat\alpha)-\tilde\beta\rVert^2)$ is $o_p(\lVert\hat\beta(\hat\alpha)-\tilde\beta\rVert)$, so it is absorbed to give $\lVert\hat\beta(\hat\alpha)-\tilde\beta\rVert=O_p(n^{-1/2})$ and hence $\lVert\hat\beta(\hat\alpha)-\tilde\beta\rVert^2=o_p(n^{-1/2})$. With the integrable second-derivative envelope $M_i$ of (A1), the full second-order remainder, including the cross term in $\partial^2_{\beta\alpha}U$ bounded through $2ab\le a^2+b^2$, is $o_p(n^{-1/2})$. Rearranging, using that $\partial_\beta U(\tilde\beta,\alpha^*)\to-A_\beta$ is invertible with probability tending to one with inverse tending to $-A_\beta^{-1}$,
\[
\sqrt{n}\,(\hat\beta(\hat\alpha)-\tilde\beta) = A_\beta^{-1}\Big(\frac{1}{\sqrt{n}}\sum_{i=1}^n U_i(\tilde\beta,\alpha^*)\Big) + o_p(1).
\]
\end{proof}

\begin{remark}[Mapping to the raking auxiliary]\label{rem:mapping}
Let $V_i=(Y_i,\tilde X_i,Z_i,A_i)$ denote the phase-one data, let $S_i(\beta)$ be the complete-data score, and let $H_i(\alpha)=H(V_i;\alpha)$ be a phase-one auxiliary constructed from an imputation model. The nuisance vector $\alpha$ may also include the projection coefficient and any preliminary coefficient estimate used to construct $H_i$. The augmented equation associated with quadratic calibration can be written
\[
U_i^{\mathrm{aug}}(\beta,\alpha)
=\frac{R_i}{\pi_i}\{S_i(\beta)-\Gamma(\beta,\alpha)H_i(\alpha)\}
+\Gamma(\beta,\alpha)H_i(\alpha).
\]
For every fixed $\alpha$, design unbiasedness gives
\[
\mathbb E\{U_i^{\mathrm{aug}}(\beta_0,\alpha)\}
=\mathbb E\{S_i(\beta_0)\}=0,
\]
because $\mathbb E(R_i/\pi_i\mid Y_i,X_i,\tilde X_i,Z_i,A_i)=1$. The expectation is therefore constant in $\alpha$, which verifies the orthogonality condition (A2). Proposition~\ref{prop:fourth-root} applies if $H_i(\alpha)$ and $\Gamma(\beta,\alpha)$ satisfy its smoothness conditions, the fitted nuisance satisfies $\lVert\hat\alpha-\alpha^*\rVert=o_p(n^{-1/4})$, and the law-of-large-numbers and central-limit conditions in (A1) hold under the phase-two sampling design. Here $n$ is the phase-two sample size; when $n/N$ tends to a positive constant, $n$ and $N$ give the same rate up to a constant. Estimating the calibration auxiliary then leaves the first-order distribution unchanged, although its probability limit can change efficiency. Proposition~\ref{prop:raking-aipw} transfers this argument exactly to quadratic raking, and Remark~\ref{rem:raking-distance} transfers it to first order for the entropy distance used in the thesis \citep{lumley2023fourthroot}.
\end{remark}

\section{Single-fit generative imputation and the parameter-uncertainty gap}\label{app:improper-mi}

We formalise the claim of Section~\ref{sec:uncertainty_gap} that a generative imputation model which is trained once and then resamples only its stochastic input omits fitted-parameter uncertainty from the between-imputation variance.

\emph{Notation.} Write $\mathcal{D}_{\mathrm{obs}}$ for the observed data (the phase-one variables for all $N$ units together with the validated $X_i$ on the rows with $R_i=1$) and $X_{\mathrm{mis}}$ for the unobserved validated values on the rows with $R_i=0$. Let $\hat{\beta}_{\mathrm{com}}$ be the complete-data estimator of the analysis parameter and $U$ its complete-data variance estimator. An imputation model $P(X_{\mathrm{mis}}\mid\mathcal{D}_{\mathrm{obs}};\theta)$ is indexed by a parameter $\theta$, which for a generative imputation model is the vector of trained network weights, with Bayesian posterior $P(\theta\mid\mathcal{D}_{\mathrm{obs}})$. From $m$ completed datasets Rubin's rules form $\bar{\beta}_m=m^{-1}\sum_d\hat{\beta}^{(d)}_{\mathrm{com}}$, $\bar{U}_m=m^{-1}\sum_d U^{(d)}$, the between-imputation variance $B_m=(m-1)^{-1}\sum_d(\hat{\beta}^{(d)}_{\mathrm{com}}-\bar{\beta}_m)^2$, and the total $T_m=\bar{U}_m+(1+m^{-1})B_m$ \citep{Rubin1987}. \emph{Proper} imputation draws $X_{\mathrm{mis}}^{(d)}$ from the posterior predictive $\int P(X_{\mathrm{mis}}\mid\mathcal{D}_{\mathrm{obs}};\theta)\,P(\theta\mid\mathcal{D}_{\mathrm{obs}})\,d\theta$, equivalently by drawing $\theta^{(d)}\sim P(\theta\mid\mathcal{D}_{\mathrm{obs}})$ and then $X_{\mathrm{mis}}^{(d)}\sim P(\cdot\mid\mathcal{D}_{\mathrm{obs}};\theta^{(d)})$. \emph{Single-fit} imputation fixes one trained value $\hat{\theta}$ and draws $X_{\mathrm{mis}}^{(d)}\sim P(\cdot\mid\mathcal{D}_{\mathrm{obs}};\hat{\theta})$ for every $d$, varying only the stochastic input.

\begin{proposition}[Single-fit imputation omits the parameter-uncertainty component]\label{prop:improper-mi}
Suppose that (C) the imputation and analysis models are congenial \citep{meng1994multiple}, so that proper multiple imputation is calibrated; (M) $\hat{\beta}_{\mathrm{com}}$ has finite second moment conditional on $\mathcal{D}_{\mathrm{obs}}$; and (R) both the conditional, or data-noise, variance $\theta\mapsto\mathrm{Var}(\hat{\beta}_{\mathrm{com}}\mid\mathcal{D}_{\mathrm{obs}},\theta)$ and the expected within-imputation variance $\theta\mapsto E[U\mid\mathcal{D}_{\mathrm{obs}},\theta]$ are constant over the posterior support of $\theta$, with $\hat{\theta}$ lying in that support. Then, as $m\to\infty$,
\[
B_m^{\mathrm{prop}}\xrightarrow{p}\mathrm{Var}\bigl(\hat{\beta}_{\mathrm{com}}\mid\mathcal{D}_{\mathrm{obs}}\bigr),
\qquad
B_m^{\mathrm{sf}}\xrightarrow{p}\mathrm{Var}\bigl(\hat{\beta}_{\mathrm{com}}\mid\mathcal{D}_{\mathrm{obs}},\hat{\theta}\bigr),
\]
and the two limits differ by exactly the parameter-uncertainty component,
\[
\lim_{m\to\infty}\bigl(B_m^{\mathrm{prop}}-B_m^{\mathrm{sf}}\bigr)
=\mathrm{Var}_{\theta\mid\mathcal{D}_{\mathrm{obs}}}\!\bigl(E[\hat{\beta}_{\mathrm{com}}\mid\mathcal{D}_{\mathrm{obs}},\theta]\bigr)\;\ge\;0 .
\]
Under (R), the within-imputation limits also agree, so $T_m^{\mathrm{prop}}-T_m^{\mathrm{sf}}$ converges to the displayed non-negative component. For every finite $m\ge2$, their conditional expectations differ by $(1+m^{-1})$ times that component. This comparison is between the two Rubin variance estimators formed from the same observed data; a coverage conclusion for the single-fit coefficient estimator additionally requires the two procedures to have the same first-order centre and sampling variance.
\end{proposition}

\begin{proof}
Condition throughout on $\mathcal{D}_{\mathrm{obs}}$. Under proper imputation the completed-data estimates $\hat{\beta}^{(d)}_{\mathrm{com}}=\hat{\beta}_{\mathrm{com}}(\mathcal{D}_{\mathrm{obs}},X_{\mathrm{mis}}^{(d)})$ are independent and identically distributed draws of $\hat{\beta}_{\mathrm{com}}$ given $\mathcal{D}_{\mathrm{obs}}$, with mean $E[\hat{\beta}_{\mathrm{com}}\mid\mathcal{D}_{\mathrm{obs}}]$ and variance $\mathrm{Var}(\hat{\beta}_{\mathrm{com}}\mid\mathcal{D}_{\mathrm{obs}})$; by the law of large numbers, finite under (M), the sample variance $B_m^{\mathrm{prop}}$ converges in probability to that variance. Under single-fit imputation the draws are independent and identically distributed given $(\mathcal{D}_{\mathrm{obs}},\hat{\theta})$, so $B_m^{\mathrm{sf}}\to\mathrm{Var}(\hat{\beta}_{\mathrm{com}}\mid\mathcal{D}_{\mathrm{obs}},\hat{\theta})$ by the same argument, the conditional second moment at $\hat{\theta}$ being finite because $E_{\theta\mid\mathcal{D}_{\mathrm{obs}}}[\mathrm{Var}(\hat{\beta}_{\mathrm{com}}\mid\mathcal{D}_{\mathrm{obs}},\theta)]\le\mathrm{Var}(\hat{\beta}_{\mathrm{com}}\mid\mathcal{D}_{\mathrm{obs}})<\infty$ forces the $(R)$-constant conditional variance to be finite. The law of total variance, applied to the posterior law of $\theta$ given $\mathcal{D}_{\mathrm{obs}}$, yields
\[
\mathrm{Var}\bigl(\hat{\beta}_{\mathrm{com}}\mid\mathcal{D}_{\mathrm{obs}}\bigr)
=E_{\theta\mid\mathcal{D}_{\mathrm{obs}}}\!\bigl[\mathrm{Var}(\hat{\beta}_{\mathrm{com}}\mid\mathcal{D}_{\mathrm{obs}},\theta)\bigr]
+\mathrm{Var}_{\theta\mid\mathcal{D}_{\mathrm{obs}}}\!\bigl(E[\hat{\beta}_{\mathrm{com}}\mid\mathcal{D}_{\mathrm{obs}},\theta]\bigr).
\]
By (R) the integrand of the first term is constant over the posterior support and equal to its value at $\hat{\theta}$, so $E_{\theta\mid\mathcal{D}_{\mathrm{obs}}}[\mathrm{Var}(\hat{\beta}_{\mathrm{com}}\mid\mathcal{D}_{\mathrm{obs}},\theta)]=\mathrm{Var}(\hat{\beta}_{\mathrm{com}}\mid\mathcal{D}_{\mathrm{obs}},\hat{\theta})$. Subtracting the two limiting between-variances leaves the parameter-uncertainty term, which is non-negative as a variance. The between- and within-imputation quantities are conditionally unbiased for every $m\ge2$: $E[B_m^{\mathrm{prop}}\mid\mathcal{D}_{\mathrm{obs}}]=\mathrm{Var}(\hat{\beta}_{\mathrm{com}}\mid\mathcal{D}_{\mathrm{obs}})$ and $E[B_m^{\mathrm{sf}}\mid\mathcal{D}_{\mathrm{obs}}]=\mathrm{Var}(\hat{\beta}_{\mathrm{com}}\mid\mathcal{D}_{\mathrm{obs}},\hat{\theta})$, a sample variance of independent and identically distributed draws being unbiased, while $E[\bar{U}_m^{\mathrm{prop}}\mid\mathcal{D}_{\mathrm{obs}}]=E_{\theta\mid\mathcal{D}_{\mathrm{obs}}}[E(U\mid\mathcal{D}_{\mathrm{obs}},\theta)]$ and $E[\bar{U}_m^{\mathrm{sf}}\mid\mathcal{D}_{\mathrm{obs}}]=E(U\mid\mathcal{D}_{\mathrm{obs}},\hat{\theta})$. By (R) the two within-variance terms coincide, so, using the law of total variance above and taking the conditional expectation of $T_m=\bar{U}_m+(1+m^{-1})B_m$,
\[
E[T_m^{\mathrm{prop}}\mid\mathcal{D}_{\mathrm{obs}}]-E[T_m^{\mathrm{sf}}\mid\mathcal{D}_{\mathrm{obs}}]=\bigl(1+m^{-1}\bigr)\,\mathrm{Var}_{\theta\mid\mathcal{D}_{\mathrm{obs}}}\!\bigl(E[\hat{\beta}_{\mathrm{com}}\mid\mathcal{D}_{\mathrm{obs}},\theta]\bigr)\ge0
\]
holds exactly for every $m\ge2$, an identity rather than a limit; averaging over $\mathcal{D}_{\mathrm{obs}}$ preserves the sign. Taking $m\to\infty$ also gives the asserted limit for $T_m^{\mathrm{prop}}-T_m^{\mathrm{sf}}$. Under (C), $T_m^{\mathrm{prop}}$ is calibrated for the proper-imputation estimator \citep{Rubin1987,meng1994multiple}. The coefficient estimate from single-fit imputation generally converges to a different conditional mean, so the variance gap relative to proper imputation does not alone identify its repeated-sampling variance or coverage.
\end{proof}

\begin{proposition}[Coverage consequence under first-order equivalence]\label{prop:improper-mi-coverage}
Let $\beta_0$ be scalar, let the sample size $n\to\infty$, and let the number of imputations $m=m_n\to\infty$. Suppose, in addition to the conditions of Proposition~\ref{prop:improper-mi}, that the proper and single-fit coefficient estimators are first-order equivalent,
\[
\sqrt n\,(\bar\beta_m^{\mathrm{sf}}-\bar\beta_m^{\mathrm{prop}})=o_p(1),
\]
that $\sqrt n(\bar\beta_m^{\mathrm{prop}}-\beta_0)\Rightarrow N(0,V)$ and $nT_m^{\mathrm{prop}}\xrightarrow{p}V$, and that
\[
n\{T_m^{\mathrm{prop}}-T_m^{\mathrm{sf}}\}\xrightarrow{p}\Delta
\qquad\text{for some }0<\Delta<V.
\]
Then $nT_m^{\mathrm{sf}}\xrightarrow{p}V-\Delta$, and the nominal normal-theory interval based on $T_m^{\mathrm{sf}}$ has limiting coverage below its nominal level.
\end{proposition}

\begin{proof}
First-order equivalence gives $\sqrt n(\bar\beta_m^{\mathrm{sf}}-\beta_0)\Rightarrow N(0,V)$. The two variance limits give $nT_m^{\mathrm{sf}}\xrightarrow{p}V-\Delta<V$, so Slutsky's theorem yields an interval whose limiting standardisation uses a denominator smaller than the sampling standard deviation.
\end{proof}

\begin{remark}\label{rem:improper-mi}
The single-fit scheme is Rubin's \emph{improper} imputation in the congenial Bayesian formulation: missing values are drawn at a fixed parameter rather than after drawing that parameter from its posterior \citep{robins2000inference}. For a deep generative model, $\hat\theta$ is the trained network weights and the stochastic input may be a variational latent draw, adversarial noise, or a diffusion initial state. A frequentist treatment of an estimated $\hat\theta$ has additional estimation terms that may have either sign, which the posterior-support formulation does not cover. Condition (R) isolates the parameter-uncertainty channel by holding the conditional spread and expected complete-data variance fixed over $\theta$. Constant residual variance in an imputation model is not sufficient by itself, because the completed-data estimator can still have a conditional variance that depends on the imputed design matrix. Condition (R) does hold, for example, when the completed-data estimator is an affine function of the imputed values with coefficients fixed by the observed data and $\theta$ changes only their conditional mean. A smaller validation sample can increase the omitted component by increasing posterior uncertainty in $\theta$, but the size of the component also depends on how the completed-data estimate changes with $\theta$. The small between-imputation variances observed for the single-fit procedure in Chapter~\ref{chap:SICG} are consistent with this mechanism; the simulation coverage results provide the separate evidence about undercoverage.
\end{remark}

\section{Conditional-mean imputation and the collapse of the between-imputation variance}\label{app:mean-collapse}

The previous result holds even when each completed dataset is a genuine draw from $P(X_{\mathrm{mis}}\mid\mathcal{D}_{\mathrm{obs}};\hat{\theta})$. A second, more severe degeneracy arises when the imputation procedure converges to the conditional mean, so that even the data-noise variance is lost. Let $C_i$ denote the conditioning data for unit $i$ and $X$ the target variable being imputed.

\begin{lemma}[Conditional-mean imputation collapses the between-imputation variance]\label{lemma:mean-collapse}
Suppose $E[X^2]<\infty$, the number of imputed units and the number of imputations $m$ are fixed and finite, and, in the limit indexed for instance by an increasing reconstruction-penalty weight or by training towards the population optimum, an imputation procedure produces values $\hat{X}_i^{(d)}$, $d=1,\dots,m$, whose within-unit conditional variance vanishes, $\mathrm{Var}(\hat{X}_i^{(d)}\mid C_i)\to0$, so that the imputations concentrate on a deterministic function $\hat{g}(C_i)$ of the conditioning data; suppose further that $\hat{\beta}_{\mathrm{com}}$ is continuous in the imputed values. Then $B_m\xrightarrow{p}0$ and Rubin's total variance reduces to $\bar{U}_m$. Rubin's formula then contains no between-imputation contribution, whereas that contribution is positive under a proper procedure whenever the completed-data estimates vary across posterior predictive draws. For a subsequent sequence with sample size $n\to\infty$, suppose that the collapsed estimator is first-order equivalent to a proper estimator with $\sqrt n(\bar\beta_m^{\mathrm{prop}}-\beta_0)\Rightarrow N(0,V)$, while $n\bar U_m\xrightarrow{p}V_U$ for some $0<V_U<V$. The nominal normal-theory interval based on $\bar U_m$ then undercovers. A squared-error reconstruction penalty motivates the deterministic limit, because the minimiser over square-integrable $g$ of $E[(g(C)-X)^2]$ is the conditional mean $g^{*}(C)=E[X\mid C]$.
\end{lemma}

\begin{proof}
For any square-integrable measurable $g$, so that the cross term below is integrable by the Cauchy--Schwarz inequality,
\[
E\bigl[(g(C)-X)^2\bigr]=E\bigl[(g(C)-E[X\mid C])^2\bigr]+E\bigl[\mathrm{Var}(X\mid C)\bigr],
\]
    since the cross term $E\{(g(C)-E[X\mid C])(E[X\mid C]-X)\}$ vanishes after conditioning on $C$. The first term is non-negative and is minimised, uniquely outside a null set, at $g(C)=E[X\mid C]=:g^{*}(C)$, which is deterministic in $C$. This population result motivates the vanishing-conditional-variance limit assumed in the lemma. Extending it to a fitted procedure at a finite penalty would require a uniform-convergence result for $\hat{g}$.

Taking the stated concentration on $\hat{g}(C_i)$ as given, componentwise convergence of the finite imputed vector implies joint convergence. Continuity of $\hat{\beta}_{\mathrm{com}}$ and the continuous mapping theorem give $\hat{\beta}^{(d)}_{\mathrm{com}}\xrightarrow{p}\hat{\beta}_{\mathrm{com}}(\mathcal{D}_{\mathrm{obs}},\hat{g}(C))$ for every $d$, so $B_m\xrightarrow{p}0$ and $T_m-\bar{U}_m=(1+m^{-1})B_m\xrightarrow{p}0$. Under a proper congenial procedure, the between component is positive when completed-data estimates vary across posterior predictive draws \citep{meng1994multiple}. The collapse removes that component from Rubin's variance. The lemma's additional first-order equivalence and variance-calibration conditions convert this variance gap into repeated-sampling undercoverage.
\end{proof}

\begin{remark}\label{rem:mean-collapse}
Reconstruction loss is one route to this degeneracy; mode collapse, to which adversarial training is prone \citep{yoon2018gain}, and prediction-averaging, as in \texttt{missForest} \citep{stekhoven2012missforest}, are others, since each can render the imputations nearly deterministic in $C$. The lemma is the variance-side counterpart of the finding that squared-error-minimising imputations target $E[X\mid C]$ and can score well on pointwise criteria without supplying valid multiple-imputation variation \citep{naf2023imputation}. It is distinct from Proposition~\ref{prop:improper-mi}: there the between-imputation variance retains the data-noise term and omits only the parameter-uncertainty term, whereas here the data-noise term itself is lost.
\end{remark}
